\chapter{Critical categories and cyclic realization}
\label{mc:chap:categories}

The potential of Chapter~\ref{vc:chap:critical} has three
incidence branches when the displacement is nonzero.
Each branch yields a matrix factorization.
To compare this category with a geometric coefficient,
we must retain the morphisms between the objects and
prove that the chosen objects generate the category.
The return products remain nonzero before the displacement
is inverted. After inversion, every complex of morphisms
between distinct branches contracts.

We construct all nine morphism complexes and their products
in a finite Clifford model. Multiplier derivatives then
prove generation for the polynomial family. The same
morphism comparison survives arbitrary base change,
but generation does not. A monomial example records its
multiplicity in a higher operation. Absolute Hochschild
and cyclic calculations show which of these data survive
passage to global periodic invariants.

Let $S$ be a commutative ring in even parity and $W\in S$.
A matrix factorization consists of a finite projective
$\mathbb Z/2$-graded $S$-module and an odd endomorphism $D$
satisfying $D^2=W1$. For two such objects,
\[
 \partial_{ab}(c)=D_bc-(-1)^{|c|}cD_a
\]
is the differential on their morphisms. It squares to zero.

Composition of endomorphisms satisfies the graded Leibniz
rule. All explicit factorizations below are finite free.
All categorical gradings in this chapter are parity
gradings unless an integer Koszul degree is specified.

Write $\operatorname{MF}(S,W)^\pi$ for the idempotent
completion of this dg category. A collection
\emph{split-generates} if finite sums, shifts, cones and
direct summands produce every object of its homotopy
category. The closure under these operations is its
\emph{thick closure}. For a dg algebra $A$,
$\operatorname{Perf}(A)$ denotes perfect right dg modules:
the thick closure of the free module $A$.
A Morita equivalence is an equivalence on perfect modules
with their dg enhancements.

\section{A relative Clifford comparison}

Let $R$ be a commutative $\mathbb Q$-algebra and
$A=(A_{ij})$ a symmetric $n$-by-$n$ matrix over $R$.
Put $S=R[x_1,\ldots,x_n]$.

On $S\otimes\Lambda^*R^n$, let $e_i$ be exterior
multiplication and $\iota_i$ contraction. Thus
\[
 [e_i,e_j]_+=[\iota_i,\iota_j]_+=0,
 \qquad [e_i,\iota_j]_+=\delta_{ij},
\]
where $[u,v]_+=uv+vu$ for odd operators.
Define
\begin{equation}\label{mc:tcrm:eq:quadratic-object}
 D=\sum_i x_i e_i+\sum_i(Ax)_i\iota_i,
 \qquad
 \gamma_i=e_i-\sum_jA_{ij}\iota_j.
\end{equation}
Then
\begin{equation}\label{mc:tcrm:eq:operator-relations}
 D^2=x^tAx,\quad [D,\gamma_i]_+=0,\quad
 [D,\iota_i]_+=x_i,\quad
 [\gamma_i,\gamma_j]_+=-2A_{ij}.
\end{equation}
Let $\operatorname{Cl}_R(-A)$ be the parity-graded
$R$-algebra on odd generators $\gamma_i$ with the last
relations and zero differential.

\begin{theorem}\label{mc:tcrm:thm:Clifford}
The indicated operators define a quasi-isomorphism
of $R$-linear dg algebras
\begin{equation}\label{mc:tcrm:eq:Clifford-map}
 \operatorname{Cl}_R(-A)\longrightarrow
 \bigl(\operatorname{End}_S(S\otimes\Lambda^*R^n),[D,-]\bigr).
\end{equation}
It commutes with arbitrary commutative base change of
$R$. Neither nondegeneracy of $A$ nor reducedness of
$R$ is required.
\end{theorem}

\begin{proof}
The exterior-contraction operators $e_I\iota_J$, with
each index set ordered increasingly, form an $S$-basis
of the endomorphisms. Replacing each $e_i$ by
$\gamma_i=e_i-\sum_jA_{ij}\iota_j$ is triangular in
exterior length, with diagonal entries one.
Consequently the $\gamma_I\iota_J$ are another basis.
In particular, the ordered $\gamma_I$ are linearly
independent and realize the stated Clifford algebra.

In this basis the differential kills $\gamma_I$ and
replaces each $\iota_i$ by $x_i$, with the Koszul
sign past $\gamma_I$. The underlying complex is
therefore $\operatorname{Cl}_R(-A)$ tensored with
the polynomial Koszul complex of $x_1,\ldots,x_n$.
That sequence is regular over every coefficient ring.
Its augmentation evaluates at $x=0$ and removes
positive contraction length.

For polynomial degree $d$ and contraction length $q$,
the Koszul homotopy away from $d+q=0$ is
\[
 \frac{\sum_i\iota_i\partial_{x_i}}{d+q}.
\]
On the Clifford factor it carries the sign $(-1)^{|\gamma_I|}$.
Its anticommutator with the differential is the
identity away from degree $d+q=0$.
This proves the quasi-isomorphism and its compatibility
with base change. The inclusion is multiplicative.
The projection and homotopy are only linear maps.
\end{proof}

Over a field with a nondegenerate form, this is the
usual Clifford description of quadratic matrix
factorizations.\footnote{T.~Dyckerhoff,
\emph{Compact generators in categories of matrix
factorizations},
\href{https://arxiv.org/abs/0904.4713}{arXiv:0904.4713},
Section~5.5, following the quadratic theory of
Buchweitz--Eisenbud--Herzog. The isolated-singularity
generation theorem is not needed for the endomorphism
comparison proved here.}
The theorem computes the chosen object. Generation
of a larger nonisolated category is a separate question.

The coefficient hypothesis matters for this presentation
and contraction. Over $\mathbb Z$, the anticommutator
relation with $n=1$ and $A=0$ gives
$\mathbb Z[g]/(2g^2)$, whereas the exterior operator
satisfies $g^2=0$. The element $g^2$ in that presented
algebra is nonzero, as its image in $\mathbb F_2[g]$
shows. Characteristic zero alone does not justify
either division by two or the displayed Euler homotopy.

For the central triangle take
\[
 R=\mathbb C[P,Q,T],\qquad
 A=\begin{pmatrix}P&T/2\\T/2&Q\end{pmatrix}.
\]
Then
\begin{equation}\label{mc:tcrm:eq:binary-Clifford}
 \gamma_1^2=-P,\qquad \gamma_2^2=-Q,
 \qquad \gamma_1\gamma_2+\gamma_2\gamma_1=-T.
\end{equation}
The even element $w=2\gamma_1\gamma_2+T$ satisfies
\begin{equation}\label{mc:tcrm:eq:even-cover}
 w^2=T^2-4PQ=\Delta,
 \qquad w\gamma_i=-\gamma_iw.
\end{equation}
Since the even basis is $1,\gamma_1\gamma_2$, its
subalgebra is exactly $R[w]/(w^2-\Delta)$, the
coordinate ring of the cover in \eqref{vc:cqm:eq:covers}.

It is not the center of the full universal Clifford
algebra. On the locus $\Delta\ne0$, adjoining
$\sqrt\Delta$ gives idempotents
$\tfrac12(1\pm w/\sqrt\Delta)$, exchanged by the
deck involution. This recovers its sign local system.

It does not supply a functor from matrix factorizations
to perverse sheaves.

\section{The three objects and every morphism}

Now put $R=\mathbb C[h,P,Q,T]$ and retain the
actual displaced potential
\begin{equation}\label{mc:tcrm:eq:potential}
 W_h=Px(x+h)+Qy(y+h)+Txy.
\end{equation}
Let $C_R$ be the free rank-four Clifford algebra
\eqref{mc:tcrm:eq:binary-Clifford} over this ring.
Give it the odd derivation
$\partial_0\gamma_1=hP$, $\partial_0\gamma_2=hQ$.
Set
\[
 t_0=0,\qquad t_1=h\gamma_1,\qquad t_2=h\gamma_2.
\]
Each satisfies $\partial_0t_a+t_a^2=0$.

\begin{definition}\label{mc:tcrm:def:category}
The marked dg category $\mathcal C_h$ has objects
$0,1,2$. For every ordered pair, put
\begin{equation}\label{mc:tcrm:eq:Hom}
 \operatorname{Hom}_{\mathcal C_h}(a,b)=C_R,
 \qquad
 \partial_{ab}(c)=\partial_0c+t_bc-(-1)^{|c|}ct_a.
\end{equation}
Composition from
$\operatorname{Hom}(b,c)\otimes_R\operatorname{Hom}(a,b)$
to $\operatorname{Hom}(a,c)$ is Clifford multiplication
in that order. The units are the copies of $1$.
\end{definition}

The displayed identities for $t_a$ give $\partial_{ab}^2=0$.
Expanding \eqref{mc:tcrm:eq:Hom} gives the Leibniz rule
for every composable pair. Thus the definition gives
a strict dg category, finite free in every morphism
complex.

On $R[x,y]\otimes\Lambda^*R^2$, use the Clifford
operators from \eqref{mc:tcrm:eq:quadratic-object} and put
\begin{align}
 D_0={}&xe_1+ye_2+(Px+Ty/2+hP)\iota_1
                       +(Tx/2+Qy+hQ)\iota_2,
                       \label{mc:tcrm:eq:objects}\\
 D_1={}&D_0+h\gamma_1,\qquad
 D_2=D_0+h\gamma_2.\nonumber
\end{align}
All three squares are $W_h1$. They are Koszul
factorizations at the respective sections
$(0,0)$, $(-h,0)$ and $(0,-h)$ of the incidence locus.

A dg functor is a \emph{quasi-equivalence} if it induces
quasi-isomorphisms on morphism complexes and is essentially
surjective on the homotopy category.

\begin{theorem}\label{mc:tcrm:thm:marked-model}
The constant-in-$x,y$ Clifford operators give a
quasi-equivalence from $\mathcal C_h$ to the full
dg subcategory of matrix factorizations of $W_h$
on $D_0,D_1,D_2$. It preserves composition and units
strictly and is compatible with arbitrary commutative
base change of $R$.
\end{theorem}

\begin{proof}
Write $a_0=(0,0)$, $a_1=(1,0)$ and $a_2=(0,1)$.
In every matrix morphism complex use the basis
$c\iota_J$, with $c$ a Clifford word.
The differential on $c$ is \eqref{mc:tcrm:eq:Hom}, and
for a single contraction generator it is
\begin{equation}\label{mc:tcrm:eq:source-shift}
 \partial(c\iota_i)=\partial_{ab}(c)\iota_i
       +(-1)^{|c|}c\bigl(x_i+h(a_a)_i\bigr),
 \qquad (x_1,x_2)=(x,y).
\end{equation}
Indeed $[D_a,\iota_i]_+=x_i+h(a_a)_i$.

The regular sequence belongs to the source object $a$.
Iteration gives the tensor product with its polynomial
Koszul complex. Translating coordinates reduces its
contraction to the one in
Theorem~\ref{mc:tcrm:thm:Clifford}.

Thus inclusion is a quasi-isomorphism on every morphism
complex. It preserves multiplication, and the three
objects already exhaust the stated target subcategory.
The contractions are defined over $R$, so they survive
base change, including $h=0$ and nilpotent parameters.
\end{proof}

This also permits the base change
$R\to\mathbb C[P,Q,T][[h]]$, with $x,y$ still polynomial
variables. No analytic convergence or generation of
the entire matrix-factorization category follows from
the morphism comparison alone. The generation theorem
proved below supplies the latter assertion
while the multiplier variables remain in the ring.
Arbitrary specialization can destroy it.

\section{Cohomology before and after localization}

\begin{proposition}\label{mc:tcrm:prop:off-diagonal}
For $a\ne b$ there is an odd $R$-linear operator
$H_{ab}$ on the morphism complex such that
\begin{equation}\label{mc:tcrm:eq:h-homotopy}
 \partial_{ab}H_{ab}+H_{ab}\partial_{ab}=h1.
\end{equation}
Moreover
\begin{equation}\label{mc:tcrm:eq:off-diagonal}
 H^{\bar0}\operatorname{Hom}(a,b)\cong R/(h),
 \qquad H^{\bar1}\operatorname{Hom}(a,b)\cong R/(h).
\end{equation}
\end{proposition}

\begin{proof}
Write $\delta_{ab}=\partial_{ab}/h$.
The odd Clifford derivatives $\partial_{\gamma_i}$
are defined by $\partial_{\gamma_i}\gamma_j=\delta_{ij}$
and the super Leibniz rule. They preserve the Clifford
relations, and
\[
 \delta_{ab}\partial_{\gamma_i}
       +\partial_{\gamma_i}\delta_{ab}
          =((a_b)_i-(a_a)_i)1.
\]
For distinct objects one coordinate difference is
$\pm1$. Divide by that sign to define $H_{ab}$.
Thus $\delta_{ab}$ is split exact.

In the bases $1,\gamma_1\gamma_2$ and
$\gamma_1,\gamma_2$, each of its two parity matrices
has determinant zero and an entry $\pm1$.
Its kernel and image are therefore the same free
rank-one summand. Multiplying the differential by
the non-zero-divisor $h$ keeps that kernel and
multiplies its image by $h$. This proves
\eqref{mc:tcrm:eq:off-diagonal}.
\end{proof}

For the three diagonal complexes define, respectively,
\begin{equation}\label{mc:tcrm:eq:regular-pairs}
 (r_1,r_2)=(P,Q),\quad(-P,Q-T),\quad(P-T,-Q),
 \qquad \kappa_a=r_1\gamma_2-r_2\gamma_1.
\end{equation}

\begin{proposition}\label{mc:tcrm:prop:diagonal}
For each object $a$,
\begin{equation}\label{mc:tcrm:eq:diagonal}
 H^{\bar0}\operatorname{End}(a)=R/(hr_1,hr_2),
 \qquad
 H^{\bar1}\operatorname{End}(a)=(R/(h))\kappa_a.
\end{equation}
Their odd generators satisfy
\begin{equation}\label{mc:tcrm:eq:odd-square}
 \kappa_a^2=PQ(T-P-Q)
\end{equation}
in the even cohomology ring.
\end{proposition}

\begin{proof}
The differentials of $\gamma_1$, $\gamma_2$ and
$v=\gamma_1\gamma_2$ are $hr_1$, $hr_2$ and
$h\kappa_a$. Each pair $(r_1,r_2)$ is regular.
Hence every odd cycle is a multiple of its Koszul
syzygy $\kappa_a$, and the odd boundaries are its
$h$-multiples. An even cycle is scalar. Its boundaries
form $(hr_1,hr_2)$. Finally
\[
 \kappa_a^2=-Qr_1^2-Pr_2^2+Tr_1r_2,
\]
which equals \eqref{mc:tcrm:eq:odd-square} for all three
pairs in \eqref{mc:tcrm:eq:regular-pairs}.
\end{proof}

After inverting $h$, all off-diagonal morphism
complexes contract by $H_{ab}/h$. The diagonal
cohomology rings become $R[h^{-1}]/(r_1,r_2)$,
with zero odd cohomology. These are the three
multiplier lines of Proposition~\ref{vc:cicp:prop:triangle-strata}.
The localized category is quasi-equivalent to the
disjoint union of its diagonal dg algebras. This
does not infer formality of those algebras from
their cohomology rings.

At $h=0$, by contrast, every differential in
\eqref{mc:tcrm:eq:Hom} vanishes. The endomorphism algebra
of the direct sum of the three identical objects is
$\operatorname{Mat}_3(C_R/(h))$.
Each off-diagonal derived fiber has two directions in each parity.
One comes from cohomology modulo $h$, the other from
$h$-torsion in the opposite parity.

For object zero, the odd
specialization retains the syzygy sequence
\[
 0\longrightarrow R_0\kappa_0
 \longrightarrow R_0\gamma_1\oplus R_0\gamma_2
 \xrightarrow{(a,b)\mapsto Pa+Qb}(P,Q)
 \longrightarrow0,\qquad R_0=\mathbb C[P,Q,T].
\]
The additional even vector $v$ has connecting class
$\partial v/h=\kappa_0$. Cohomology followed by ordinary
specialization would lose these connecting maps.

\section{Return and transit products}

We now fix generators, so every composition has a
specified sign. With $v=\gamma_1\gamma_2$, take
the following closed representatives:
\begin{equation}\label{mc:tcrm:eq:generators}
\begin{array}{c|c|c}
 ab&\alpha_{ab}\ (\text{even})&\beta_{ab}\ (\text{odd})\\ \hline
 01&Q+v&\gamma_1\\
 10&T-Q+v&\gamma_1\\
 02&T-P+v&\gamma_2\\
 20&P+v&\gamma_2\\
 12&T+v&\gamma_2-\gamma_1\\
 21&v&\gamma_2-\gamma_1.
\end{array}
\end{equation}
They generate the two groups in
\eqref{mc:tcrm:eq:off-diagonal}. Put $K=P+Q-T$ and
$\ell_{01}=P$, $\ell_{02}=Q$, $\ell_{12}=K$.
The notation $\ell_{ab}$ is symmetric in its indices.

\begin{proposition}\label{mc:tcrm:prop:composites}
For $a\ne b$, the return products in
$H^*\operatorname{End}(a)$ are
\begin{align}
 [\beta_{ba}\beta_{ab}]&=-\ell_{ab},\nonumber\\
 [\alpha_{ba}\alpha_{ab}]&=-
       \prod_{\{i,j\}\ne\{a,b\}}\ell_{ij},
       \label{mc:tcrm:eq:returns}\\
 [\alpha_{ba}\beta_{ab}]&=\varepsilon_{ab}\kappa_a,
 \qquad
 [\beta_{ba}\alpha_{ab}]=-\varepsilon_{ab}\kappa_a.
       \nonumber
\end{align}
The product ranges over the other two edges.
Set $\varepsilon_{ab}=1$ for $a<b$ and $-1$ for $a>b$.
All these classes are nonzero and killed by $h$.

For distinct $a,b,c$, write the four transit products as
\[
 \begin{aligned}
 \alpha_{bc}\alpha_{ab}&=u_1\alpha_{ac},&
 \beta_{bc}\beta_{ab}&=u_2\alpha_{ac},\\
 \alpha_{bc}\beta_{ab}&=u_3\beta_{ac},&
 \beta_{bc}\alpha_{ab}&=u_4\beta_{ac}.
 \end{aligned}
\]
Their coefficients are
\begin{equation}\label{mc:tcrm:eq:transit}
\begin{array}{c|rrrr}
 abc&u_1&u_2&u_3&u_4\\ \hline
 012&Q&-1&P&K\\
 021&-P&-1&-Q&K\\
 102&-K&-1&P&-Q\\
 120&P&1&-K&Q\\
 201&K&1&Q&-P\\
 210&-Q&1&-K&-P.
\end{array}
\end{equation}
These identities hold already for the chosen chain
representatives.
\end{proposition}

\begin{proof}
The relations \eqref{mc:tcrm:eq:binary-Clifford} give
\[
 v^2=-Tv-PQ,\qquad
 v\gamma_1=P\gamma_2-T\gamma_1,\qquad
 \gamma_1v=-P\gamma_2,\qquad
 v\gamma_2=-Q\gamma_1,\qquad
 \gamma_2v=Q\gamma_1-T\gamma_2.
\]
Substituting \eqref{mc:tcrm:eq:generators} proves
\eqref{mc:tcrm:eq:returns} and every entry of
\eqref{mc:tcrm:eq:transit}. For example, the path
$0\to1\to0$ gives
\[
 -P,\quad -QK,\quad \kappa_0,\quad-\kappa_0
\]
in the order odd--odd, even--even, even--odd,
odd--even. The transit product
$(\gamma_2-\gamma_1)\gamma_1=P-T-v=-\alpha_{02}$
gives the second entry in the first transit row.

Every boundary has a factor $h$, whereas each
displayed return is a nonzero polynomial Clifford
vector modulo $h$. Thus none is a boundary.
Equation~\eqref{mc:tcrm:eq:h-homotopy} supplies a primitive
for $h$ times either off-diagonal factor. Multiplying
that primitive by the other closed factor, with its
Leibniz sign, proves that each return is killed by $h$.
\end{proof}

The three plane forms have a symmetric expression
in the original constraint multipliers:
\begin{equation}\label{mc:tcrm:eq:symmetric}
 P=p+r,\qquad Q=q+r,\qquad T=2r,
 \qquad
 \begin{aligned}
 \Delta&=-4(pq+qr+rp),\\
 -\kappa_a^2&=(p+q)(q+r)(r+p).
 \end{aligned}
\end{equation}
Indeed the potential becomes
\[
 p\,x(x+h)+q\,y(y+h)+r(x+y)(x+y+h).
\]
The transformations $(x,y)\mapsto(y,x)$ with
$p\leftrightarrow q$, and
$(x,y)\mapsto(-(x+y+h),y)$ with
$p\leftrightarrow r$, preserve it. They permute
the three branch points and the three plane forms.
This is the $S_3$ symmetry of the potential and its
marked support. An equivariant categorical comparison
must additionally specify the maps intertwining its
chosen factorizations.

\section{Generation from smooth branches}

Let $S$ be a regular Noetherian ring of finite Krull
dimension, let $W$ be a non-zero-divisor, and set
$A=S/(W)$. The singularity category is
$D_{\mathrm{sg}}(A)=D^b_{\mathrm{coh}}(A)/\operatorname{Perf}(A)$.
The regular-ambient equivalence identifies it with
matrix factorizations. The image of an $A$-module under
this equivalence is its \emph{stabilization}.\footnote{
D.~Orlov, \emph{Matrix factorizations for nonaffine LG-models},
\href{https://arxiv.org/abs/1101.4051}{arXiv:1101.4051},
Theorem~3.5. An affine regular Noetherian scheme of finite
Krull dimension has the required locally free resolutions.

For the local and nonlocal ring formulations, see
T.~Dyckerhoff, \emph{Compact generators in categories of
matrix factorizations}.
\href{https://arxiv.org/abs/0904.4713}{arXiv:0904.4713},
Sections~2.1--2.4.}

\begin{proposition}\label{mc:cgcr:prop:branches}
Let $f_1,\ldots,f_r\in S$ act null-homotopically on
every finite matrix factorization, and put $I=(f_1,\ldots,f_r)$.
Suppose ideals $J_1,\ldots,J_k$ contain $W$, satisfy
$J_1\cdots J_k\subset I$, and have regular quotients $S/J_a$.
Then the stabilizations of the $S/J_a$ split-generate
$\operatorname{MF}(S,W)^\pi$.
\end{proposition}

\begin{proof}
If $DH+HD=f\,1$, the cone of multiplication by $f$ has
differential and change of basis
\begin{equation}\label{mc:cgcr:eq:cone}
 D_f=\begin{pmatrix}D&f\,1\\0&-D\end{pmatrix},
 \qquad U=\begin{pmatrix}1&-H\\0&1\end{pmatrix},
 \qquad D_fU=U\begin{pmatrix}D&0\\0&-D\end{pmatrix}.
\end{equation}
The second summand is shifted, so $U$ is even and invertible.
Thus the original object is a direct summand of its cone.

Repeat for every $f_i$. On a split sum with differential
$\operatorname{diag}(D,-D)$, a subsequent homotopy is
$\operatorname{diag}(H_i,-H_i)$. Conjugate it back through
the preceding splitting. No anticommutation of the initially
chosen homotopies is needed.

Choose a bounded coherent representative $N$ in
$D_{\mathrm{sg}}(A)$. Its ordinary Koszul tensor
\[
 N\otimes K(f_1,\ldots,f_r)
\]
represents these iterated cones.
Every cohomology module $V$ is killed by $I$.

Filter $V$ by successive products of the $J_a$.
The product containment makes the last term zero, and
each quotient is a finite module over one regular ring
$S/J_a$. It therefore has a finite projective resolution
over that ring. Each projective term is a summand of a
finite sum of $S/J_a$ as an $A$-module.

The filtration and the bounded truncation triangles put
the Koszul tensor in the thick closure of the branch
modules. Its original summand belongs to that closure
after idempotent completion.
\end{proof}

The derivatives of $W$ supply the homotopies. For a free
factorization,
\begin{equation}\label{mc:cgcr:eq:Jacobian-homotopy}
 D(\partial_zD)+(\partial_zD)D=(\partial_zW)1.
\end{equation}
For a projective module, choose a connection along the
derivation $\partial_z$ and differentiate the endomorphism
$D$ covariantly. Such a connection exists by writing the
module as an idempotent summand of a free module and
projecting its coefficientwise connection. The same
Leibniz identity proves \eqref{mc:cgcr:eq:Jacobian-homotopy}.
Flatness of the connection is unnecessary.

\subsection{The polynomial triangle}

Put $R=\mathbb C[h,P,Q,T]$, $S=R[x,y]$, and retain
\[
 W_h=Px(x+h)+Qy(y+h)+Txy.
\]
Let $D_0,D_1,D_2$ be the factorizations
\eqref{mc:tcrm:eq:objects}, and let $\mathcal C_h$ be the
finite dg category of Definition~\ref{mc:tcrm:def:category}.
Write $\mathcal E_h=\bigoplus_{a,b=0}^2
\operatorname{Hom}_{\mathcal C_h}(a,b)$, with matrix
composition as multiplication. Its underlying $R$-module
has rank $36$.

\begin{theorem}\label{mc:cgcr:thm:triangle}
The object $G=D_0\oplus D_1\oplus D_2$ split-generates
$\operatorname{MF}(S,W_h)^\pi$. There is a Morita equivalence
\begin{equation}\label{mc:cgcr:eq:Morita}
 \operatorname{Perf}(\mathcal E_h)
       \simeq\operatorname{MF}(S,W_h)^\pi.
\end{equation}
At $h=0$, with $P,Q,T$ still variables, $D_0|_{h=0}$
alone split-generates. Its Morita algebra is the rank-four
Clifford algebra \eqref{mc:tcrm:eq:binary-Clifford} over
$R_0=\mathbb C[P,Q,T]$, with zero differential.
\end{theorem}

\begin{proof}
The multiplier derivatives generate
\[
 I=(x(x+h),y(y+h),xy).
\]
Set $J_0=(x,y)$, $J_1=(x+h,y)$ and $J_2=(x,y+h)$.
The invertible change
\[
 z_0=h+x+y,\qquad z_1=-x,\qquad z_2=-y
\]
identifies $I$ with $(z_0z_1,z_0z_2,z_1z_2)$.
It follows that $I=J_0\cap J_1\cap J_2$ and
$J_0J_1J_2\subset I$. All three quotients are $R$.

For a module killed by $I$, the filtration
\[
 0\subset J_1J_2V\subset J_1V\subset V
\]
has quotients killed by $J_0,J_2,J_1$, respectively.
Their projective resolutions over $R$ have length at most
four. Equations~\eqref{mc:cgcr:eq:Jacobian-homotopy} and
Proposition~\ref{mc:cgcr:prop:branches} apply.

Each $D_a$ is the Koszul stabilization of its branch
module. Indeed its exterior coefficients generate $J_a$
and its contraction coefficients express $W_h$ in that
regular sequence. The Koszul stabilization is obtained
by adding the corresponding multiplication-by-$W_h$
homotopy to the Koszul resolution.\footnote{
T.~Dyckerhoff, \emph{Compact generators in categories of
matrix factorizations}.
\href{https://arxiv.org/abs/0904.4713}{arXiv:0904.4713},
Section~2.3, Lemma~2.6 and Corollary~2.7. Section~2.4
extends the construction to regular nonlocal rings.

Exterior--contraction duality changes the stabilization
by a parity shift, which leaves its thick closure unchanged.}

The full morphism comparison of
Theorem~\ref{mc:tcrm:thm:marked-model} is multiplicative.
It identifies $\mathcal E_h$ with the derived endomorphisms
of $G$. The two Morita functors are
$\operatorname{RHom}(G,-)$ and
$-\otimes^{\mathbb L}_{\mathcal E_h}G$.

At $h=0$ the derivative ideal becomes $(x,y)^2$.
The two-step $(x,y)$-adic filtration reduces every module
killed by this ideal to modules over the regular ring
$R_0$. Proposition~\ref{mc:cgcr:prop:branches}, with the branch
ideal repeated twice, proves the central assertion.
Theorem~\ref{mc:tcrm:thm:Clifford} supplies its endomorphism
algebra.
\end{proof}

The proof also applies to finite projective factorizations
over $\mathbb C[P,Q,T][[h]][x,y]$.
The branch quotients remain regular, the filtrations are
finite, and the morphism contractions are coefficient-linear.
The cyclic comparisons
below concern the finite-type algebraic models.

\section{Failure under multiplier specialization}

\begin{proposition}\label{mc:cgcr:prop:specialization}
For $R\to\mathbb C$ given by $h=P=Q=T=0$, the natural functor
\[
 \operatorname{Perf}(\mathcal E_h\otimes_R^{\mathbb L}\mathbb C)
       \longrightarrow\operatorname{MF}(\mathbb A^2,0)^\pi
\]
is not essentially surjective.
\end{proposition}

\begin{proof}
All three images are
$K_0=(\mathbb C[x,y]\otimes\Lambda^*\mathbb C^2,xe_1+ye_2)$.
On $x\ne0$, contraction by $\iota_1/x$ contracts $K_0$.
On $y\ne0$, use $\iota_2/y$. Its thick closure vanishes
away from the origin. The free even factorization
$(\mathcal O_{\mathbb A^2},0)$ does not vanish there.
It therefore lies outside the image.
\end{proof}

The coefficient-linear contractions in
Theorem~\ref{mc:tcrm:thm:marked-model} do survive this base
change, so the functor still computes the full subcategory
generated by $K_0$. What fails is generation of the new
ambient category: an arbitrary zero-potential object has
no multiplier derivatives from which to obtain the
homotopies \eqref{mc:cgcr:eq:Jacobian-homotopy}.

\section{A monomial and its higher products}
\label{mc:rmf:sec:monomial}

One object already distinguishes a higher operation from
its periodic invariant. The potential $\lambda z^m$ has
multiplicity $m$, while its binary cohomology algebra is
the same for every $m\ge3$.
We keep $\lambda$ as a parameter and compute the complete
minimal algebra of its chosen factorization.

Put $R=\mathbb C[\lambda]$, $S=R[z]$ and $m\ge2$.
All complexes in this section are $\mathbb Z/2$-graded,
with even scalar coefficients. On the free module
$S^{1|1}$ define
\begin{equation}\label{mc:rmf:eq:factorization}
 D_m=\begin{pmatrix}0&z\\\lambda z^{m-1}&0\end{pmatrix},
 \qquad D_m^2=\lambda z^m I.
\end{equation}
Let $A_m=\operatorname{End}_S(S^{1|1})$, with differential
$\delta a=D_ma-(-1)^{|a|}aD_m$.
The square of this differential is zero because $D_m^2$
is scalar.

\subsection{The parameter-linear contraction}

Set
\begin{equation}\label{mc:rmf:eq:matrices}
 J=\begin{pmatrix}0&-1\\\lambda z^{m-2}&0\end{pmatrix},
 \qquad H=\begin{pmatrix}0&0\\1&0\end{pmatrix},
 \qquad E=\begin{pmatrix}1&0\\0&0\end{pmatrix}.
\end{equation}
The first two matrices are odd and the third is even.
Direct multiplication gives
\begin{equation}\label{mc:rmf:eq:pairs}
 \delta J=0,\qquad \delta H=zI,\qquad
 \delta E=zJ,\qquad J^2=-\lambda z^{m-2}I.
\end{equation}
Every even matrix has a unique expression $aI+bE$,
and every odd matrix a unique expression $aJ+bH$,
with $a,b\in R[z]$.

\begin{proposition}\label{mc:rmf:prop:contraction}
The cohomology of $A_m$ is the free $R$-module
$R1\oplus R\chi$, with $1$ even and $\chi=[J]$ odd.
For $a_0=a|_{z=0}$, define
\begin{align*}
 p(aI+bE)&=a_0\,1,&p(aJ+bH)&=a_0\chi,\\
 i(1)&=I,&i(\chi)&=J,\\
 \mathsf h(aI+bE)&=\frac{a-a_0}{z}H,
 &\mathsf h(aJ+bH)&=\frac{a-a_0}{z}E.
\end{align*}
Then $pi=1$, $\delta\mathsf h+\mathsf h\delta=1-ip$ and
$\mathsf h^2=\mathsf h i=p\mathsf h=0$. These identities survive every
commutative base change of $R$.
\end{proposition}

\begin{proof}
Equations~\eqref{mc:rmf:eq:pairs} split the underlying
complex into two copies of multiplication by $z$:
one sends $H$ to $zI$, and the other sends $E$ to $zJ$.
This computes the cohomology. The polynomial $a-a_0$
is divisible by $z$, so the displayed homotopy uses
no localization. Applying $\delta\mathsf h+\mathsf h\delta$ to
$aI+bE$ gives $(a-a_0)I+bE$, and applying it to
$aJ+bH$ gives $(a-a_0)J+bH$.
The other identities follow from the same decomposition.
All formulas are polynomial in $\lambda$ and require
no flatness of its specialization.
\end{proof}

\subsection{Every arity from one polynomial identity}

Let $s$ denote parity suspension, so $|sa|=|a|-1$.
Our suspended differential and multiplication on a dg
algebra are
\begin{equation}\label{mc:rmf:eq:bar-signs}
 b_1(sa)=-s(\delta a),\qquad
 b_2(sa,sb)=(-1)^{|a|}s(ab).
\end{equation}
The maps $b_n$ extend by the tensor signs to a degree-one
coderivation on the tensor coalgebra of the suspension.
An $A_\infty$ structure means that this coderivation
squares to zero. A morphism means a coalgebra map
commuting with these coderivations.

Introduce an even auxiliary variable $\tau$ and put
\begin{equation}\label{mc:rmf:eq:translated}
 D_m(\tau)=
 \begin{pmatrix}
  0&z-\tau\\
  \displaystyle\lambda\sum_{j=0}^{m-1}z^{m-1-j}\tau^j&0
 \end{pmatrix}.
\end{equation}
The geometric sum gives
\[
 D_m(\tau)^2=\lambda(z^m-\tau^m)I.
\]
Writing $a(\tau)=D_m(\tau)-D_m=\sum_{j=1}^{m-1}\tau^jA_j$,
we obtain
\begin{equation}\label{mc:rmf:eq:Taylor}
 A_1=J,\qquad A_j=\lambda z^{m-1-j}H\quad(2\le j<m),
 \qquad \delta a(\tau)+a(\tau)^2=-\lambda \tau^mI.
\end{equation}

\begin{theorem}\label{mc:rmf:thm:minimal}
On $R1\oplus R\chi$, put the strict unit products and
the single additional reduced coefficient
\begin{equation}\label{mc:rmf:eq:higher}
 b_m((s\chi)^{\otimes m})=\lambda s1.
\end{equation}
A reduced coefficient has no unit among its inputs.
All other reduced coefficients vanish. These operations
give a strictly unital minimal $A_\infty$ algebra.

There is an $A_\infty$ quasi-isomorphism to $A_m$ with
\begin{equation}\label{mc:rmf:eq:comparison}
 F_1(s1)=sI,\qquad
 F_j((s\chi)^{\otimes j})=sA_j\quad(1\le j<m).
\end{equation}
All remaining components are zero, including every
$F_j$ with $j>1$ and a unit input.
The formulas commute with arbitrary commutative
base change of $R$.
\end{theorem}

\begin{proof}
Every word without units is a power of $s\chi$,
which has even parity. On such words the equation
$b_{A_m}F=Fb$ has generating series
\[
 -s\bigl(\delta a(\tau)+a(\tau)^2\bigr)
       =\lambda \tau^m sI.
\]
This is \eqref{mc:rmf:eq:Taylor}. On the right, a source
$b_m$ creates $s1$. A higher component of $F$ with that
unit input is zero. Only $F_1(s1)$ at arity $m$ remains.

For words containing a unit, strict unitality fixes
$b_2(s1,sa)=sa$ and
$b_2(sa,s1)=(-1)^{|a|}sa$.
Deleting a unit on its left or on its right gives
opposite terms after the tensor sign. Internal terms
therefore cancel in pairs. At either endpoint the
remaining term equals multiplication by the target
unit $I$. Every higher coefficient with a unit input
vanishes. This proves the morphism equation on all words.

The same cancellation proves $b^2=0$. For $m\ge3$,
an inner $b_m$ creates a unit, so it cannot be an input
to another normalized higher coefficient. The only
remaining terms are the two adjacent binary unit terms,
which cancel. For $m=2$, the assertion is associativity
of the algebra with unit and relation $\chi^2=-\lambda$.
Finally $F_1=i$ induces the cohomology isomorphism of
Proposition~\ref{mc:rmf:prop:contraction}. The polynomial
formulas and that contraction persist after base change.
\end{proof}

For $m=2$, \eqref{mc:rmf:eq:higher} is the Clifford
product $\chi^2=-\lambda$. For $m\ge3$, the binary
product has $\chi^2=0$, while the first nonzero reduced
higher operation has arity $m$ when $\lambda\ne0$.
At $\lambda=0$, every reduced coefficient vanishes.
The model becomes the exterior algebra on $\chi$.

The field version relating higher multiplications to
potential coefficients is Dyckerhoff's theorem.\footnote{
T.~Dyckerhoff, \emph{Compact generators in categories of
matrix factorizations}.
\href{https://arxiv.org/abs/0904.4713}{arXiv:0904.4713},
Section~5.6, Theorem~5.8. The present formulas specify
the coefficient parameter, the suspension signs and
the complete comparison for the chosen object.}

\begin{corollary}\label{mc:rmf:cor:multiplicity}
Fix a nonzero complex value of $\lambda$ and $m\ge3$.
Among these strictly unital minimal algebras, the arity
of the first nonzero reduced higher operation is
preserved by $A_\infty$ isomorphisms. In particular,
it determines $m$.
\end{corollary}

\begin{proof}
A parity-preserving unital isomorphism has linear term
$\chi\mapsto c\chi$ with $c\ne0$.
On a pure $s\chi$ word, every Taylor component has even
parity in the suspension. Its value therefore lies in
the line spanned by $s\chi$.

Before the first higher operation, the morphism
equations involve only products of two odd elements,
which are zero. At that first arity, the nonzero
$s1$ coefficient can only be matched by a higher
operation of the same arity in the other algebra.
Its coefficient changes by the nonzero factor $c^m$.
The first arity is therefore unchanged.
\end{proof}

\begin{proposition}\label{mc:cgcr:prop:monomial}
With $\lambda$ retained as a variable, $D_m$ split-generates
$\operatorname{MF}(\mathbb C[\lambda,z],\lambda z^m)^\pi$.
\end{proposition}

\begin{proof}
The multiplier derivative of any factorization differential
contracts $z^m$ by \eqref{mc:cgcr:eq:Jacobian-homotopy}.
Use $m$ copies of the branch ideal $(z)$ in
Proposition~\ref{mc:cgcr:prop:branches}. Their product is
$(z^m)$ and the quotient is the regular ring
$\mathbb C[\lambda]$. Its Koszul stabilization is $D_m$.
\end{proof}

Thus the minimal algebra above gives a Morita model of
the total polynomial family. After specializing $\lambda$,
its explicit endomorphism comparison remains valid.
That comparison alone does not prove generation of the
specialized ambient category.

\section{Absolute Hochschild coefficients}

All invariants in this section are over $\mathbb C$.
For a complex dg category $\mathcal A$, its Hochschild
cochain complex is the derived endomorphism complex of
the diagonal bimodule:
\[
 C^\bullet(\mathcal A)=R\operatorname{Hom}_{
     \mathcal A^{\mathrm{op}}\otimes_{\mathbb C}\mathcal A}
                  (\mathcal A,\mathcal A).
\]
The diagonal has component $\operatorname{Hom}_{\mathcal A}(a,b)$
at the object pair $(a,b)$.
Its cohomology is $\operatorname{HH}^*(\mathcal A)$.

Suppose a smooth Landau--Ginzburg model has all critical
points in the zero fiber. Its absolute Hochschild cohomology
is computed by polyvectors with differential given by
contraction with $dW$.\footnote{K.~H.~Lin and D.~Pomerleano,
\emph{Global matrix factorizations}.
\href{https://arxiv.org/abs/1101.5847}{arXiv:1101.5847},
Theorem~3.1 and Lemma~3.3. The complex is absolute,
so it includes multiplier derivatives.}

Return to $R_0=\mathbb C[P,Q,T]$ and $S_0=R_0[x,y]$.
For $W_0=Px^2+Txy+Qy^2\in S_0$, the identity
$W_0=P\partial_PW_0+Q\partial_QW_0+T\partial_TW_0$
verifies the critical-value hypothesis.

\begin{proposition}\label{mc:cgcr:prop:Hochschild}
The absolute Hochschild cohomology of
$\operatorname{MF}(S_0,W_0)^\pi$ is the parity folding
of the cohomology of
\begin{equation}\label{mc:cgcr:eq:Jacobian}
 \mathcal J=\frac{R_0[x,y]}{(x^2,y^2)}
                 \otimes\Lambda(\alpha,\beta,\gamma),
 \quad
 \begin{cases}
 d\alpha=xy,\\
 d\beta=2Px+Ty,\\
 d\gamma=Tx+2Qy,
 \end{cases}
\end{equation}
where $|\alpha|=|\beta|=|\gamma|=-1$.
As an $R_0$-complex, its free ranks in degrees
$0,-1,-2,-3$ are $4,12,12,4$.
Moreover,
\begin{equation}\label{mc:cgcr:eq:zeroth}
 H^0(\mathcal J)=R_0\oplus M,\qquad M^2=0,
 \qquad M=\operatorname{coker}
       \begin{pmatrix}2P&T\\T&2Q\end{pmatrix}.
\end{equation}
The dimensions of its derived fibers are
\[
\begin{array}{c|rrrr}
 \text{rank of the quadratic form}&H^0&H^{-1}&H^{-2}&H^{-3}\\\hline
 2&1&3&3&1\\
 1&2&6&6&2\\
 0&3&9&9&3.
\end{array}
\]
\end{proposition}

\begin{proof}
The five gradient entries are
$x^2,y^2,xy,2Px+Ty,Tx+2Qy$.
The first two are a regular sequence. Tensoring their
Koszul resolution with the other free Koszul factors
and then augmenting gives \eqref{mc:cgcr:eq:Jacobian}.
In degree zero all products of $x,y$ vanish, while their
linear relations are the columns in \eqref{mc:cgcr:eq:zeroth}.

At a rank-two form the two linear gradients are independent
coordinates. Eliminating them leaves three closed exterior
generators over $\mathbb C$.

At a rank-one form, make a
linear change to a nonzero square. One linear gradient
and the square of the remaining coordinate form a regular
sequence, with a length-two quotient. The three excess
Koszul generators are closed after elementary generator
changes.

At the zero form, multiplication by $xy$ on
$\mathbb C[x,y]/(x^2,y^2)$ has rank one. Its kernel and
cokernel both have dimension three. The other two exterior
generators are closed. These three computations give the table.
\end{proof}

Here $H^0$ means integer Koszul degree zero, rather than
the entire even group after folding. The module $M$ is
supported on $4PQ-T^2=0$, has rank one at its generic
point, and needs two generators at the vertex.

The strict supercenter consists of elements that
supercommute with every Clifford generator.
For the universal Clifford algebra it is $R_0$.
An even central coefficient of $\gamma_1\gamma_2$ must vanish.
An odd central vector lies in the kernel of the displayed
Hessian, which is zero over the polynomial domain. At $P=Q=T=0$ the Clifford
algebra becomes exterior and is entirely supercentral.
Thus strict supercenters also need not commute with base
change.

The table describes fibers of the fixed absolute
complex. It makes no identification with invariants of
every newly specialized matrix-factorization category.

\subsection{The monomial Hochschild modules}

For $F_m=\lambda z^m$, the absolute gradient has two entries.
They retain the multiplier derivative even when one later
takes a fiber over the multiplier base.

\begin{proposition}\label{mc:prop:monomial-Hochschild}
For $m\ge2$, the absolute gradient Koszul complex has
\begin{align}
 H^0&=\mathbb C[\lambda,z]/(z^m,\lambda z^{m-1}),\notag\\
 H^{-1}&=\bigl(\mathbb C[\lambda,z]/(z^{m-1})\bigr)\theta,
 &\theta&=m\lambda\epsilon_\lambda-z\epsilon_z,
 \label{mc:cgcr:eq:monomial-HH}\\
 H^{-2}&=0.\notag
\end{align}
Its parity folding gives absolute Hochschild cohomology.
\end{proposition}

\begin{proof}
The differential is
$d\epsilon_\lambda=z^m$,
$d\epsilon_z=m\lambda z^{m-1}$.
Cancel $z^{m-1}$ in the syzygy equation in the polynomial
domain. The regular pair $(z,m\lambda)$ has syzygies
generated by $(m\lambda,-z)$.
The degree-minus-two boundary is $-z^{m-1}\theta$,
with zero kernel. The degree-zero quotient gives the
remaining assertion.
\end{proof}

\section{The mixed cyclic comparison}

For a complex dg category $\mathcal A$, let
$CC_\bullet(\mathcal A)$ be its absolute Hochschild bar
complex, with internal and Hochschild differential $b$
and Connes differential $B$. These odd operators satisfy
$b^2=B^2=bB+Bb=0$. Its periodic cyclic homology is the
cohomology of $(CC_\bullet(\mathcal A)((u)),b+uB)$,
where $u$ is an even formal variable. Laurent series have
finitely many negative powers and arbitrary nonnegative
powers of $u$. For
$\mathcal A=\mathcal C_h$, the bar words range over all
composable cycles of its three objects.

A $u$-connection is an operator $\nabla_u=\partial_u+A(u)$,
where $A(u)$ is $\mathbb C((u))$-linear and
$[\nabla_u,b+uB]=(b+uB)/(2u)$.
A mixed-complex map $f$ is weakly compatible with
connections if there is an odd $\mathbb C((u))$-linear map
$H$ such that
\[
 \nabla'_u f-f\nabla_u=(b'+uB')H+H(b+uB).
\]
Here $f$ also denotes its extension to Laurent series.

Return to $R=\mathbb C[h,P,Q,T]$ and $S=R[x,y]$.
Efimov's comparison, applied to the Morita equivalence
\eqref{mc:cgcr:eq:Morita}, gives a zigzag of quasi-isomorphisms
of mixed complexes
\begin{equation}\label{mc:cgcr:eq:cyclic}
 (CC_\bullet(\mathcal E_h),b,B)
     \simeq(\Omega_S^\bullet,-dW_h\wedge,d).
\end{equation}
The comparison is weakly compatible with the connections.

It passes from curved modules through the completed
curved bar complexes for $(S,-W_h)$ and $(S,W_h)$.
Their completion is the product over bar-word lengths.
It allows the curvature term to insert an additional
bar entry while retaining an infinite sequence of lengths.

The final curved Hochschild--Kostant--Rosenberg map is
\[
 a_0[a_1|\cdots|a_k]\longmapsto
       \frac{a_0}{k!}\,da_1\wedge\cdots\wedge da_k.
\]
Its inputs are scalar curved-bar words after those
comparisons, not Clifford words.\footnote{A.~I.~Efimov,
\emph{Cyclic homology of categories of matrix factorizations}.
\href{https://arxiv.org/abs/1212.2859v2}{arXiv:1212.2859v2},
Theorems~1.1--1.4, Proposition~3.14 and Section~4,
equations~(4.2)--(4.6). The hypotheses are smoothness,
separation, finite type and characteristic zero, with
the critical locus contained in the zero fiber.
The linear multiplier identity verifies the last condition.}

On the de Rham side the differential and connection are
\begin{equation}\label{mc:cgcr:eq:connection}
 Q_u=-dW_h\wedge+ud,\qquad
 \nabla_u=\partial_u+\frac{\Gamma}{u}+\frac{W_h}{u^2},
 \qquad \Gamma|_{\Omega^p}=-p/2.
\end{equation}
They satisfy $Q_u^2=0$ and
$[\nabla_u,Q_u]=Q_u/(2u)$.
Indeed the four nonzero commutator terms are
$d$, $(dW_h\wedge)/(2u)$, $-d/2$ and
$-(dW_h\wedge)/u$. Their sum is the stated expression.
All six affine coordinates, including $h,P,Q,T$, are
differentiated. For $h$ fixed, the corresponding
five-variable absolute model is a separate application.

Even a length-one bar boundary uses both return paths.
In the convention where it is the supercommutator,
closed odd maps $f:0\to1$ and $g:1\to0$, both represented
by $\gamma_1$, give
\[
 b(g[f])=gf-(-1)^{|f||g|}fg=-P1_0-P1_1.
\]
This relates the two returns. It does not kill each term
separately or permit the intermediate object to be deleted.

\section{The grading of the periodic realization}

Let $M$ be a smooth quasi-projective complex $d$-fold
with potential $F$.
Form the quotient of its dg category of coherent matrix
factorizations by the thick subcategory generated by
totalizations of short exact sequences.
Write $\operatorname{MF}(M,F)$ for its idempotent completion.
Assume that every critical point lies in $F^{-1}(0)$.

Write $\phi_F$ for classical vanishing cycles and normalize
$\Phi_F=\phi_F[-1]\mathbb Q_M[d]$.
We regard this complex on $M$ by extension by zero.
Efimov's periodic comparison identifies
$\operatorname{HP}_k$ with the parity folding of
$H^{k-1}(F^{-1}(0),\phi_F\mathbb C_M)((u))$,
with connection monodromy $T_F(-1)^k$.

Consequently
\begin{equation}\label{mc:cgcr:eq:degrees}
 \begin{aligned}
 H^j(M,\Phi_F)&\longleftrightarrow
 H^{j+d-1}(F^{-1}(0),\phi_F\mathbb Q_M),\\
 k&=j+d.
 \end{aligned}
\end{equation}
Here $T_F$ is monodromy in the value of $F$.
The support convention is ordinary cohomology. The multiplier
projection need not be proper.

For the Morse function $xy$ on $\mathbb A^2$, the Milnor
fiber is homotopy equivalent to a circle.
Its reduced cohomology is in degree one, giving even
periodic degree $k=2$.
For $x^2$ on $\mathbb A^1$, the two-point Milnor
fiber has reduced cohomology in degree zero, giving odd
periodic degree $k=1$. These examples fix the shift.

For the central triangle,
Proposition~\ref{vc:cqm:prop:integral} gives
$R\Gamma(\mathbb A^3,\Phi_{W_0})=\mathbb Q[1]$.
The single normalized degree is $j=-1$ and $d=5$.
Thus the raw vanishing degree is $3$ and the periodic
degree is $4$. We obtain
\begin{equation}\label{mc:cgcr:eq:central-periodic}
 \operatorname{HP}_{\bar0}\operatorname{MF}(\mathbb A^5,W_0)
      =\mathbb C((u)),\qquad
 \operatorname{HP}_{\bar1}=0.
\end{equation}

The nonzero summand $\mathcal P_-$ of
Theorem~\ref{vc:cqm:thm:coefficient} is self-dual.
Its compact cohomology vanishes by
Proposition~\ref{vc:cqm:prop:integral}, so Verdier duality
also makes its ordinary cohomology zero.
It contributes nothing to these global periodic groups.

Raw function monodromy is $+1$ on this class. The even
periodic degree leaves that sign unchanged.

At a fixed
nonzero numerical $h$, the three multiplier lines of
Proposition~\ref{vc:cicp:prop:triangle-strata} are Morse--Bott
with four hyperbolic normal coordinates. Their orientation
local systems are constant on the affine lines. Each
contributes one even class, giving periodic ranks $(3,0)$.
These are separately specified absolute categories.

\begin{corollary}\label{mc:cor:monomial-periodic}
For every $m\ge2$, the absolute periodic cyclic ranks
of $\operatorname{MF}(\mathbb A^2,\lambda z^m)^\pi$
are $(1,0)$.
\end{corollary}

\begin{proof}
For $\Phi_m=\varphi_{F_m}\mathbb Q[2]$, dimensional
reduction over the multiplier has
$R\pi_!\Phi_m=\mathbb Q_{\{z=0\}}$ in rational sheaves.
Verdier duality gives the same ordinary pushforward.
Its normalized ordinary cohomology therefore has one
class in degree zero. The ambient dimension is two,
so \eqref{mc:cgcr:eq:degrees} puts it in even periodic degree.
\end{proof}

The absolute Hochschild modules in
\eqref{mc:cgcr:eq:monomial-HH} depend on $m$, as does the
first reduced operation in Theorem~\ref{mc:rmf:thm:minimal}.
The constant periodic ranks determine neither datum.
No assertion about conservativity of a categorical functor
follows from this equality of ranks alone.

\section{Global elliptic sectors and periodic ranks}

Let $E$ be a complex elliptic curve,
$\Theta=\mathcal O_E(0)$ its theta line, and $\vartheta$
its divisor section. Fix $n\ge1$, set $h=[n]t$, and put
$B_n=E_x\times E_y\times E_t$.
For $u_1=x,u_2=y,u_3=x+y$, define
\[
 \mathcal S_i=u_i^*\Theta\otimes(u_i+[n]t)^*\Theta,
 \qquad s_i=\vartheta(u_i)\vartheta(u_i+[n]t).
\]
Let $V_n=\bigoplus_i\mathcal S_i$,
$M_n=\operatorname{Tot}_{B_n}(V_n^\vee)$, and
$F_n=\langle\lambda,s\rangle$.
This is a smooth quasi-projective sixfold with a global
regular function. Fiberwise multiplier dilation shows
that every critical point lies in its zero fiber.

\begin{proposition}\label{mc:cgcr:prop:extra-sectors}
Above each point with $[n]t\in E[2]\setminus\{0\}$ and
$x=y=-[n]t$, the potential at $\lambda=0$ has analytic
local form $pu+qv+rw$. The factorization
\begin{equation}\label{mc:cgcr:eq:extra}
 D_{\mathrm{extra}}=ue_1+ve_2+we_3
                         +p\iota_1+q\iota_2+r\iota_3
\end{equation}
on $\mathcal O\otimes\Lambda^*\mathbb C^3$ is a generator over the convergent local ring of this
chart, and over its formal completion. Its endomorphism
cohomology is $\mathbb C$ in
even parity and zero in odd parity. The three moving
branch objects vanish near these points. There are
$3n^2$ such local sectors.
\end{proposition}

\begin{proof}
The fourth branch and its torsion condition are those of
Theorem~\ref{ei:esf:thm:global-fibers} and
\eqref{ei:esf:eq:fourth-branch}.
The first two sections have simple roots at the indicated
fourth branch point. Use them as coordinates $u,v$.
Subtract their multiples from the third section. On
their common branch its restriction is
$\vartheta(-2h)\vartheta(-h)$.

At nonzero two-torsion this has a simple zero in $h$,
since $[2]$ is \'{e}tale and $\vartheta(-h)\ne0$.
It is therefore a third coordinate $w$.
The dual change of multiplier frame gives $F_n=pu+qv+rw$.
The point of the $t$-curve is also unramified under $[n]$.

Exterior--contraction relations prove
$D_{\mathrm{extra}}^2=F_n1$.
The endomorphism differential sends
$\iota_i$ to $(u,v,w)_i$ and $e_i$ to $(p,q,r)_i$.
The resulting six-coordinate Koszul complex resolves
$\mathbb C$. Generation follows by the smooth-branch
argument with $(u,v,w)$, or by iterated hyperbolic
stabilization. The three other branch supports are
disjoint from this point, so their objects restrict to
zero there. Finally each of the three nonzero points
of $E[2]$ has $n^2$ preimages under $[n]$.
\end{proof}

At fixed $w=0$ the relative chart has two hyperbolic
pairs and a free multiplier direction $r$. The proposition
uses the absolute chart, in which $w$ is a coordinate.
It constructs local generators. A global finite generator
with all descent maps is a further construction.

Write $Z_n$ for the reduced common zero locus of the $s_i$.
Its curve part consists of three copies of $E$,
\[
 C_0=(0,0,t),\qquad C_1=(0,-[n]t,t),\qquad
 C_2=(-[n]t,0,t).
\]
They meet at the $n^2$ points $t\in E[n]$.
The remaining points are exactly the $3n^2$ points of
Proposition~\ref{mc:cgcr:prop:extra-sectors}.

\begin{theorem}\label{mc:cgcr:thm:global-periodic}
The Betti numbers and absolute periodic cyclic ranks are
\begin{align}
 (b_0,b_1,b_2)(Z_n)&=(1+3n^2,\,2n^2+4,\,3),
       \label{mc:cgcr:eq:Betti}\\
 \dim_{\mathbb C((u))}\operatorname{HP}_{\bar0}
             \operatorname{MF}(M_n,F_n)&=3n^2+4,\notag\\
 \dim_{\mathbb C((u))}\operatorname{HP}_{\bar1}
             \operatorname{MF}(M_n,F_n)&=2n^2+4.
       \label{mc:cgcr:eq:global-periodic}
\end{align}
\end{theorem}

\begin{proof}
The first two constraints force $x,y\in\{0,-[n]t\}$.
Evaluating the third gives exactly the stated locus.
For $C=C_0\cup C_1\cup C_2$, normalization has the
exact sheaf sequence
\begin{equation}\label{mc:cgcr:eq:normalization}
 0\longrightarrow\mathbb Q_C
 \longrightarrow\bigoplus_{i=0}^2(\nu_i)_*\mathbb Q_E
 \longrightarrow\bigoplus_{t\in E[n]}\mathbb Q_t^2
 \longrightarrow0.
\end{equation}
The last map is $(a,b,c)\mapsto(b-a,c-a)$ at every
intersection. Its map on global constants has rank two.
The long exact sequence therefore gives
$b_0(C)=1$, $b_1(C)=6+(2n^2-2)$ and $b_2(C)=3$.

The isolated points give \eqref{mc:cgcr:eq:Betti}.
In particular $\chi(Z_n)=n^2$.

Let $\pi:M_n\to B_n$ and $i:Z_n\hookrightarrow B_n$.
For $\Phi_n=\phi_{F_n}[-1]\mathbb Q_{M_n}[6]$,
linear-potential dimensional reduction gives
\begin{equation}\label{mc:cgcr:eq:reduction}
 R\pi_!\Phi_n\simeq i_*\mathbb Q_{Z_n}.
\end{equation}
The base dimension and vector-bundle rank are both three.

This formula is for rational sheaves. Its mixed-Hodge
version includes the Tate twist $(-3)$.\footnote{
B.~Davison, \emph{The critical CoHA of a quiver with
potential}, \href{https://arxiv.org/abs/1311.7172}{arXiv:1311.7172},
Theorem~A.1. T.~Kinjo, \emph{Dimensional reduction in
cohomological Donaldson--Thomas theory},
\href{https://arxiv.org/abs/2102.01568}{arXiv:2102.01568},
equation~(1.1) and Theorem~4.14.
Use the critical-chart orientation and its comparison
with the canonical shifted-cotangent orientation.}

Verdier self-duality of $\Phi_n$ turns the compact
pushforward into
$R\pi_*\Phi_n\simeq i_*\omega_{Z_n}$,
where $\omega_{Z_n}$ is the rational dualizing complex.
Since $Z_n$ is proper,
\[
 H^j(M_n,\Phi_n)=H^{-j}(Z_n,\mathbb Q)^\vee.
\]

The normalized degrees are $0,-1,-2$.
Equation~\eqref{mc:cgcr:eq:degrees} gives raw vanishing
degrees $5,4,3$ and periodic degrees $6,5,4$.
Thus the even periodic rank is $b_0+b_2$ and the odd
rank is $b_1$, proving \eqref{mc:cgcr:eq:global-periodic}.
\end{proof}

The ranks are $(7,6)$ on the displacement base $n=1$
and $(16,12)$ on the half-displacement cover $n=2$.
They are invariants of the absolute global spectral
category. A relative cyclic complex over the $t$-curve,
its base-change maps and its connection in $t$ require
their own construction. The periodic $u$-connection has
the comparison above. The rank formula does not compute
its full monodromy matrix.

\section{The coherent comparison and its grading}

The same equations have the coherent Koszul algebra
\begin{equation}\label{mc:tcrm:eq:coherent-Koszul}
 K_h=\mathbb C[h,x,y]\otimes\Lambda(\epsilon_1,\epsilon_2,\epsilon_3),
 \qquad
 \begin{cases}
 d\epsilon_1=x(x+h),\\
 d\epsilon_2=y(y+h),\\
 d\epsilon_3=xy,
 \end{cases}
 \qquad |\epsilon_i|=-1.
\end{equation}
Its coherent derived category is related to an
equivariant singularity category of the potential,
with the multiplier grading retained.
Precisely, let $X=\mathbb A^3_{h,x,y}$ and let
$Z=\{W_h=0\}\subset X\times\mathbb A^3_{P,Q,T}$.
The multiplicative group dilates $P,Q,T$ with
weight one. The equivariant singularity category
is the quotient of bounded equivariant coherent
complexes on $Z$ by equivariant perfect complexes.

Isik's Koszul duality identifies it with the bounded
coherent derived category of $K_h$. The nonregular
section requires the Koszul dg algebra in place of
its degree-zero cohomology.\footnote{M.~U.~Isik,
\emph{Equivalence of the derived category of a variety
with a singularity category}.
\href{https://arxiv.org/abs/1011.1484}{arXiv:1011.1484},
Theorem~3.6 and Remark~3.7. Here the smooth ambient
scheme is the absolute space $\mathbb A^3_{h,x,y}$.
The section is $(x(x+h),y(y+h),xy)$.}

This graded equivalence has a different source of
grading from the parity category in
Definition~\ref{mc:tcrm:def:category}. To use it for
the marked objects, one must specify their graded
lifts, the parameter-linear functor and its effect
on the objects and their compositions.

\begin{question}\label{mc:tcrm:qn:coherent-comparison}
Construct the graded lifts and the $\mathbb C[h]$-linear
Koszul-duality comparison.
Require it to carry the point modules at $(0,0)$, $(-h,0)$
and $(0,-h)$ of \eqref{mc:tcrm:eq:coherent-Koszul} to the
objects \eqref{mc:tcrm:eq:objects}.
Specify the shifts and periodization. How does it identify their
Yoneda products with \eqref{mc:tcrm:eq:returns} and
\eqref{mc:tcrm:eq:transit}, and the coherent higher
operations with the relative operations of the
marked category?
\end{question}

\section{The compact categorical comparison}

The Morita equivalence \eqref{mc:cgcr:eq:Morita} identifies
two descriptions of the specified polynomial potential.
An identification with a compact sheaf deformation
potential requires an additional geometric map.
The half-displacement cover, the quadratic determinant
cover and the Enriques canonical cover have different bases.
Comparing their lines requires maps between those bases.
The genus-two parameter family and weight-five automorphic
line also remain separate constructions.

\begin{question}\label{mc:cgcr:qn:compact-comparison}
Construct the comparison from compact translation-reduced
critical Hall coefficients to the spectral category or
to an independently specified state endomorphism category.
Which maps preserve extension pullback and pushforward,
orientation, the common translation torsor, the relative
empty object and the reduced state? How do they carry
the return products of \eqref{mc:tcrm:eq:returns}, the
unintegrated summand $\mathcal P_-$ of
Theorem~\ref{vc:cqm:thm:coefficient}, and the extra local
objects \eqref{mc:cgcr:eq:extra}?

Construct the coefficient comparison before global sections
and before inverting $h$. Can it retain the arity in
Theorem~\ref{mc:rmf:thm:minimal} and the parameter specialization
map of Theorem~\ref{vc:vic:thm:specialization}?
Give a finite global generator of the elliptic spectral
category with its descent maps, including the extra local
sectors. Identify the reduced marked state and the modular
line and section required by the Igusa denominator.
\end{question}
